\documentclass[10pt,a4paper]{amsart}
\usepackage[margin=19mm]{geometry}
\usepackage{amsmath,amssymb,mathtools}
\usepackage[scr=rsfs,cal=euler]{mathalfa}
\usepackage{xfrac}
\usepackage[unicode,hidelinks,bookmarks=false]{hyperref}
\newcommand{\ie}{i.e.\ }
\newcommand{\cf}{cf.\ }
\newcommand{\resp}{respectively\ }
\newcommand{\Subsection}{\subsection}
\providecommand{\nobreakdash}{\nobreak\hbox{-}}
\setlength{\emergencystretch}{2em}
\multlinegap0pt
\usepackage{bm}
\usepackage{bbm}
\usepackage{dsfont}
\usepackage{xspace}
\usepackage{cancel}
\usepackage{mathtools}
\usepackage{amsthm}
\makeatletter
\@ifundefined{lemma}{\newtheorem{lemma}{Lemma}}{}
\@ifundefined{remark}{\newtheorem*{remark}{Remark}}{}
\makeatother
\newcommand\numberthis{\addtocounter{equation}{1}\tag{\theequation}}
\newcommand{\ts}{\tilde{\Sigma}}
\title[Proof of the bounded conformal conjecture]{Proof of the bounded conformal conjecture}
\author{Sameer Kumar}
\date{Source: arXiv:2306.15322v4 (CC BY 4.0)}
\begin{document}
\maketitle

\noindent\emph{Source and licence.} Proof of the bounded conformal conjecture. Version arXiv:2306.15322v4 (CC BY 4.0), distributed under the Creative Commons Attribution 4.0 International licence, as recorded in the arXiv licence metadata for this exact version and shown on the abstract page https://arxiv.org/abs/2306.15322v4. Full rights notice: licenses/arxiv-2306.15322-CC-BY-4.0.txt.\par\smallskip

\noindent\textit{Editorial scope.} Counted unit: the Lemma whose source label is \texttt{lemma: slice algebra} in the article (source0.tex lines 911--922, source label \texttt{lemma: slice algebra}), delivered together with the article's own complete original proof of it (source0.tex lines 924--960, one \texttt{proof} environment of 37 lines). No mathematics is written, repaired or re-derived: statement and proof are the article's own text line for line. Required context retained uncounted: none. Every definition, display and label of those lines is retained unchanged. Omitted from this package and not redistributed: 1--910, 923--923, 961--1300. Nothing inside a retained excerpt is omitted or altered: every retained body is delivered exactly as the article's lines are written, so no omission receipt of any kind accompanies this package. Citations: the retained excerpts carry no citation command at all, so no bibliography entry of the article is transcribed and no reference is redistributed. Reference adaptation: none was needed and none was made; every reference inside the delivered unit resolves inside it, so no unresolved reference or citation is produced. Printed slips kept literal: no printed word, symbol, hypothesis or number of the counted statement or of its proof was altered. Layout and class: the article's own class is \texttt{amsart} with packages including fontenc, amsmath, amssymb, amsthm, amscd, comment, cite, amsfonts, indentfirst, color, setspace, bbold, euscript, mathtext; the wrapper is the lane's pinned \texttt{amsart} editorial wrapper at 10pt on A4, which restates the theorem-like environments the retained text needs and adds the article's own macro definitions that the retained text needs (\texttt{\textbackslash numberthis}, \texttt{\textbackslash ts}), with extra wrapper packages (bm, bbm, dsfont, xspace, cancel, mathtools). The delivered numbering is the unit's own. Assets: no retained span contains a figure environment, a graphics-inclusion command, a PDF-page inclusion, a table or a TikZ picture; no image file is copied into this package, the package declares no asset dependency and no image file is redistributed with it. Licence: version arXiv:2306.15322v4 of this article is distributed under the Creative Commons Attribution 4.0 International licence, as recorded in the arXiv licence metadata for this exact version and shown on the abstract page https://arxiv.org/abs/2306.15322v4. A full rights notice is delivered at licenses/arxiv-2306.15322-CC-BY-4.0.txt. Characters: every retained excerpt is pure ASCII, so the delivered engine, XeLaTeX with its default Latin Modern text font, reports no missing character.
\begin{lemma} \label{lemma: slice algebra}
    Let $U$ be a surface in $M$, f be a Lipschitz function on $\overline{B}^{+}$, $a \in \mathbb{R}^{+}$, $\tilde{f} = \sqrt{x^2 + y^2}$, and $\varphi_z$ be defined as above for some fixed $\alpha > 0$. Let us denote $C_{\alpha} = \int\limits_{0}^{1} (\varphi_z \circ \tilde{f}) dz$. Let $\omega$ be a 1-form defined on the set U. Let $\gamma$ be a smooth function defined on $\mathbb{R}^{+}$ such that its composition $\gamma \circ f$ generates the set $\{f < a\}$ (as above). Then as $b \uparrow a$,
    \begin{align*}
        \partial (U \llcorner (\{f < a\} \cap C_{\alpha})) (\omega) - (\partial U) \llcorner (\{f < a\} \cap C_{\alpha})(\omega) = \\ 
        = - (U \llcorner C_{\alpha})(\gamma'(f)df \wedge \omega) + \int\limits_{0}^{1} dz \Big( (U \llcorner \{ f < a \})(\omega \wedge (\varphi_{z}'(\tilde{f})d\tilde{f})) \Big)
    \end{align*}
    
\begin{remark} \nonumber
    Here the symbol $C_{\alpha}$ is used to denote the function composition $\int\limits_{0}^{1} (\varphi_z \circ \tilde{f}) dz$ without the limit $b' \uparrow \alpha z$ that needs to be taken under the integral sign. We therefore, prove this lemma for the ``broken" limit of the cone at point $p \in \ts \cap \Sigma$. We will let $b' \uparrow \alpha z$ under the integral sign in the subsequent lemma. This has been done for convenience.
\end{remark}
    
\end{lemma}

\begin{proof}

    \begin{align*}
        \partial(U \llcorner (\{f < a\} \cap C_{\alpha}))(\omega) - ((\partial U) \llcorner (\{f < a\} \cap C_{\alpha}))(\omega) = \\
        = \partial((U \llcorner C_{\alpha}) \llcorner \{f < a \})(\omega) - ((\partial U \llcorner C_{\alpha}) \llcorner \{f < a \})(\omega) \; \text{(Use set rep)} \\
        = \lim\limits_{b \uparrow a} \Big[ \left( U \llcorner C_{\alpha})(( \gamma \circ f )(d\omega)) - (\partial U \llcorner C_{\alpha})(( \gamma \circ f )(\omega)\right) \Big] \label{eqn: bound_cur} \numberthis \\
    \end{align*}
    
    where in the last line we have used the fact that $\gamma \circ f$ generates $\{f < a\}$, and the definition of boundary current. We now evaluate the second term in the expression above.
    
    \begin{gather*}
        (\partial U \llcorner C_{\alpha})((\gamma \circ f)\omega) = \int\limits_{\partial U \cap C_{\alpha}}(\gamma \circ f)\omega = \int\limits_{\partial U}\left( (\gamma \circ f)\omega \int\limits_{0}^{1}(\varphi_z \circ \tilde{f})dz  \right) \\ \stackrel{Fubini}{=\joinrel=\joinrel=} \int\limits_{0}^{1} dz \left( \int\limits_{\partial U}(\gamma \circ f)\omega (\varphi_z \circ \tilde{f})  \right)
    \end{gather*}
    
    The above chain of equations can be understood as the boundary of the surface $U$ intersecting the cone $C_{\alpha}$, for some fixed $\alpha$, being generated by unifying the slices of boundary of surface $U$, intersecting each level set of the cone at fixed $z$, over all values of $z \in [0,1]$. Continuing this chain further, we have
    
    \begin{gather*}
        \stackrel{Stokes'}{=\joinrel=\joinrel=} \int\limits_{0}^{1} dz\left( \int\limits_{U} d[(\varphi_z \circ \tilde{f})((\gamma \circ f)\omega)]) \right) \\ 
        = \int\limits_{0}^{1}dz \left[ \int\limits_{U}d((\gamma \circ f)\omega) (\varphi_z \circ \tilde{f}) + (-1)^1 \int\limits_{U}((\gamma \circ f)\omega) \wedge d(\varphi_z \circ \tilde{f}) \right] \\ 
        = \int\limits_{0}^{1} dz\int\limits_{U}d((\gamma \circ f)\omega) (\varphi_z \circ \tilde{f}) - \int\limits_{0}^{1} dz\int\limits_{U} ((\gamma \circ f)\omega) \wedge \varphi_{z}'(\tilde{f})d\tilde{f}
    \end{gather*}
    
    Putting this back into \ref{eqn: bound_cur}, we get
    
    \begin{gather*}
        \lim\limits_{b \uparrow a} \bigg( ((U \llcorner C_{\alpha}))((\gamma \circ f)d\omega) - \int\limits_{0}^{1} dz \int\limits_{U}d((\gamma \circ f)\omega)(\varphi_z \circ \tilde{f}) \\ + \int\limits_{0}^{1} dz \int\limits_{U} ((\gamma \circ f)\omega) \wedge \varphi_{z}'(\tilde{f})d\tilde{f}  
        \bigg) \\
        = \lim\limits_{b \uparrow a} \bigg( ((U \llcorner C_{\alpha}))((\gamma \circ f)d\omega) - \int\limits_{U} d((\gamma \circ f)\omega) \Big( \int\limits_{0}^{1} dz (\varphi_z \circ \tilde{f}) \Big) \\  + \int\limits_{0}^{1} dz \Big( (U \llcorner \{ f < a \}) \omega \wedge (\varphi_{z}'(\tilde{f})d\tilde{f}) \Big)
        \bigg) \\
        = \lim\limits_{b \uparrow a} \bigg( ((U \llcorner C_{\alpha}))((\gamma \circ f)d\omega) 
        - \partial (U \llcorner C_{\alpha}) (\gamma \circ f\omega) \\  
        + \int\limits_{0}^{1} dz \Big( (U \llcorner \{ f < a \}) \omega \wedge (\varphi_{z}'(\tilde{f})d\tilde{f}) \Big)
        \bigg) \\
        = - (U \llcorner C_{\alpha})(\gamma'(f)df \wedge \omega) + \int\limits_{0}^{1} dz \Big( (U \llcorner \{ f < a \}) \omega \wedge (\varphi_{z}'(\tilde{f})d\tilde{f}) \Big)
    \end{gather*}
    
\end{proof} 

\end{document}
